\chapter{Kiểm thử và kết quả}
\label{chap:verification}

\section{Chiến lược kiểm thử}

Bộ test Python bao phủ parser, config, RBAC repository, rule detection,
scoring, context profiler, reporter, application service, API và CLI. Test
integration chạy từ input dataset tới artifact; test subprocess kiểm tra exit
code và thông báo CLI không lộ traceback. Frontend dùng Vitest để kiểm tra
access control presentation, transaction/approval flow, copy sản phẩm và các
helper hiển thị. Next.js build là checkpoint bổ sung cho route và type kiểm
tra ở cấp ứng dụng.

Các lệnh chất lượng lõi là:

\begin{lstlisting}[language=bash]
uv run pytest -q --cov=rbac_guard --cov-report=term-missing --cov-fail-under=85
cd web && npm test && npm run build
\end{lstlisting}

\section{Kết quả baseline}

Dataset \texttt{data/logs\_demo.csv} chứa 60 Event hợp lệ: 30 benign, 10 SQL
Injection, 10 password guessing và 10 unauthorized access. Cấu hình dùng năm
lần đăng nhập lỗi trong 300 giây; evaluate đối chiếu event ID với nhãn kỳ vọng
sau khi analyze hoàn tất. Run metadata ghi nhận 20 alert và không có dòng lỗi.
Password guessing sinh một alert cấp chuỗi nhưng alert đó mang 10 event ID, do
đó metric vẫn đánh giá đúng mười Event tấn công.

\begin{table}[htbp]
\centering
\begin{tabular}{lrrrrrrr}
\toprule
Risk type & TP & FP & FN & TN & Precision & Recall & F1 \\
\midrule
SQL Injection & 9 & 0 & 1 & 50 & 1.0000 & 0.9000 & 0.9474 \\
Password guessing & 10 & 0 & 0 & 50 & 1.0000 & 1.0000 & 1.0000 \\
Unauthorized access & 10 & 0 & 0 & 50 & 1.0000 & 1.0000 & 1.0000 \\
Macro average & -- & -- & -- & -- & 1.0000 & 0.9667 & 0.9825 \\
\bottomrule
\end{tabular}
\caption{Kết quả baseline được lấy từ \texttt{artifacts/metrics.json}.}
\label{tab:baseline}
\end{table}

Kết quả cho thấy các rule password guessing và unauthorized access khớp hoàn
toàn với scenario đã thiết kế. SQL Injection bỏ sót một biến thể trong mười
event, minh họa trade-off của regex: bổ sung pattern có thể tăng recall nhưng
cũng có nguy cơ tạo false positive. Không diễn giải bảng này là bằng chứng về
khả năng tổng quát trên log thực tế.

Vì dataset là tập scenario tổng hợp được xây dựng để kiểm chứng các rule hiện
có, bảng~\ref{tab:baseline} nên được hiểu là kết quả \emph{regression} của
implementation, không phải benchmark độc lập. Trường \texttt{expected\_label}
không được detector sử dụng, nên không có rò nhãn trong lúc phát hiện; tuy
nhiên, tập đánh giá vẫn chưa độc lập với lựa chọn pattern, ngưỡng và kịch bản
đã thiết kế. Ba mươi Event benign cũng chưa bao phủ các trường hợp ``hard
negative'' như lỗi đăng nhập hợp lệ lặp lại, IP mới do thay đổi mạng hoặc request
có từ khóa gần với SQL nhưng không phải tấn công. Do đó, precision bằng 1.0000
chỉ xác nhận không có cảnh báo sai trong phạm vi scenario này.

Password guessing còn cho thấy khác biệt giữa đơn vị đánh giá và đơn vị cảnh
báo: mười Event dương thuộc cùng một user--IP và được gom thành một alert cấp
chuỗi. Vì vậy, mười true positive ở cấp Event không tương đương với mười chiến
dịch tấn công độc lập. Một đánh giá mở rộng cần báo cáo đồng thời metric theo
Event, alert và campaign/incident, cùng tập hold-out gồm các biến thể chưa được
dùng để hình thành rule.

\section{Kết quả context-risk}

Dataset \texttt{data/logs\_risk\_demo.csv} có 14 Event hợp lệ. Khi bật
\texttt{--context-risk}, run tạo 15 alert, 10 context finding và 3 incident.
Finding gồm sáu tín hiệu ngoài giờ, hai tài nguyên hiếm, một IP mới và một
cụm từ chối lặp lại. Một Event có thể mang nhiều finding, vì vậy chúng là bằng
chứng ưu tiên hóa triage chứ không phải mẫu độc lập để tính precision/recall.
Dataset này cũng chưa có nhãn độc lập ở cấp finding hoặc incident; các con số
hiện tại xác nhận luồng tương quan và truy vết evidence, chưa định lượng tỷ lệ
phát hiện đúng hay cảnh báo sai của heuristic ngữ cảnh.

Incident nhiều tín hiệu nhất gắn với \texttt{teller01} từ IP
\texttt{198.51.100.50} trong 22:00--22:05. Nó liên kết sáu Event với mười
alert: đăng nhập lỗi, IP mới, hoạt động ngoài giờ, tài nguyên hiếm và thao tác
\texttt{users:delete} bị RBAC từ chối. Điểm lớn nhất cộng chain bonus tạo 99,
tương ứng Critical. Case này cho thấy cách hệ thống giữ quan hệ từ incident về
alert và Event gốc để nhà phân tích có thể kiểm chứng evidence.

\section{Tái lập số liệu}

Số liệu không được nhập tay. Lệnh analyze sinh artifact, sau đó evaluate sinh
metric; cùng input và cấu hình tạo alert ID ổn định từ rule ID và event ID. Các
lệnh đầy đủ nằm ở Phụ lục~\ref{app:reproduction}. Khi cần lấy số liệu mới,
nên ghi ra thư mục output riêng thay vì chỉnh sửa các artifact mẫu trong repo.
